% =====================================================================
%  6. ANÁLISIS DE HALLAZGOS  <- LA SECCIÓN QUE MÁS CALIFICA
%  Responsable: INTEGRANTE 4
%  ESTADO: REDACTADO (1 de agosto de 2026) — ~3.700 palabras + matriz completa
%  Fuentes: HRIA + CAO/IFC + estudio de salud PHR/Michigan + monitoreo COPAE
%
%  LA MATRIZ (Tabla 6) ESTÁ COMPLETA, con diez hallazgos sobre los siete
%  ejes de la auditoría. Va en página apaisada y usa longtable.
%
%  DECISIONES METODOLÓGICAS QUE CONVIENE NO DESHACER:
%
%  1) El eje ambiental se DESDOBLA en tres hallazgos distintos: la
%     insuficiencia del monitoreo y de la línea base (que sí es
%     demostrable), la afectación efectiva del agua (que NO lo es, porque
%     la evidencia está en conflicto) y la planificación del cierre.
%     Meterlos en una sola fila obligaría a clasificar como incumplimiento
%     algo que no está probado, o a rebajar algo que sí lo está.
%
%  2) La columna «Estado de cierre» dice «no verificable» en la mayoría de
%     las filas. NO es una laguna del trabajo: es el hallazgo de la
%     sección 6.7. La empresa publicó actualizaciones de su plan de acción
%     en 2010, 2011 y 2017, pero no consta verificación independiente del
%     cierre de cada recomendación.
%
%  3) Se usa el condicional al describir el hallazgo laboral («habría
%     desalentado») porque así lo formula la fuente. No lo conviertas en
%     afirmación categórica sin ir al informe original.
% =====================================================================

\section{Análisis de hallazgos}
\label{sec:hallazgos}

Esta sección sistematiza los resultados de la auditoría analizada y los convierte
en hallazgos clasificados. Se organiza en tres movimientos: primero se declara la
regla de clasificación empleada; después se presenta la matriz completa de
hallazgos sobre los siete ejes evaluados; y finalmente se desarrollan en prosa los
cuatro hallazgos de mayor gravedad, aquellos que explican por sí solos la
trayectoria del caso.

% ---------------------------------------------------------------------
\subsection{Criterio de clasificación de los hallazgos}

Antes de clasificar es obligado explicitar con qué regla se clasifica. Un auditor
que aplica categorías sin haberlas definido convierte su juicio en
inverificable, que es exactamente lo contrario de lo que persigue el principio de
enfoque basado en la evidencia.

Debe advertirse, en primer lugar, que la evaluación analizada \textbf{no clasificó
sus resultados como no conformidades}. No era una auditoría de sistema de gestión
sino una evaluación de impactos en derechos humanos, según se estableció en la
sección \ref{sec:auditoria}, y sus resultados se formularon como hallazgos y
recomendaciones por ejes temáticos. La clasificación que se presenta a
continuación es, por tanto, una \textbf{adaptación propia} construida sobre las
definiciones de la ISO 19011:2018 expuestas en la sección \ref{sec:marco}
\parencite{iso2018iso19011}. Reconocerlo no debilita el análisis: lo hace
auditable.

La regla aplicada es la siguiente:

\begin{itemize}
  \item \NCmayor: incumplimiento de un requisito que afecta derechos
        fundamentales, o que tiene carácter sistémico por comprometer la
        capacidad del conjunto del sistema de gestión para producir resultados
        conformes.
  \item \NCmenor: incumplimiento puntual, acotado y subsanable, que no
        compromete la integridad del sistema.
  \item \OBS: situación en la que no existe incumplimiento demostrado, pero se
        identifica un riesgo relevante o una oportunidad de mejora. Se emplea
        también cuando la evidencia disponible es contradictoria y no permite
        concluir.
\end{itemize}

La tercera categoría requiere una precisión, porque su uso en este caso es
deliberado y no evasivo. Cuando dos cuerpos de evidencia técnica apuntan en
direcciones opuestas y ninguno puede descartarse, el auditor no está autorizado a
elegir el que confirme su hipótesis. Registrar una observación es, en esa
situación, la única conclusión que la evidencia sostiene.

Finalmente, se entiende por \textbf{evidencia objetiva}, en un análisis de
naturaleza documental como el presente, cualquiera de estas tres fuentes: el
contenido del propio informe de auditoría, un dato de monitoreo ambiental o
sanitario procedente de un estudio identificable, o el pronunciamiento formal de
un órgano oficial nacional o internacional.

% ---------------------------------------------------------------------
\subsection{Matriz de hallazgos y no conformidades}

La Tabla \ref{tab:matriz-nc} presenta los diez hallazgos identificados sobre los
siete ejes evaluados. Su distribución es la siguiente: cinco no conformidades
mayores, cuatro menores y una observación.

Conviene señalar por qué el eje ambiental aparece desdoblado en tres filas. La
insuficiencia del monitoreo y de la línea base es un hecho documentado y
constituye por sí misma un incumplimiento. La afectación efectiva de la calidad
del agua, en cambio, es objeto de evidencia contradictoria y no puede afirmarse
como incumplimiento. Y la planificación del cierre presenta un estado distinto de
ambas. Agrupar las tres cuestiones en un único hallazgo obligaría a dar por
probado lo que no lo está, o a rebajar lo que sí lo está.

\begin{landscape}
\begin{footnotesize}
\setlength{\tabcolsep}{4pt}
\begin{longtable}{@{}L{1.7cm} L{4.2cm} L{4.6cm} L{2.9cm} L{1.6cm} L{4.4cm} L{3.2cm}@{}}
  \caption{Matriz de hallazgos y no conformidades de la auditoría socioambiental de la Mina Marlin}
  \label{tab:matriz-nc} \\
  \toprule
  \textbf{Eje} & \textbf{Hallazgo} & \textbf{Evidencia objetiva} &
  \textbf{Criterio incumplido} & \textbf{Clasif.} & \textbf{Recomendación} &
  \textbf{Estado de cierre} \\
  \midrule
  \endfirsthead
  \multicolumn{7}{@{}l}{\textit{Tabla \ref{tab:matriz-nc} (continuación)}} \\[2pt]
  \toprule
  \textbf{Eje} & \textbf{Hallazgo} & \textbf{Evidencia objetiva} &
  \textbf{Criterio incumplido} & \textbf{Clasif.} & \textbf{Recomendación} &
  \textbf{Estado de cierre} \\
  \midrule
  \endhead
  \bottomrule
  \multicolumn{7}{@{}p{23cm}}{\textit{Nota.} Elaboración propia a partir de
  \textcite{oncommonground2010hria}, \textcite{cao2006marlin},
  \textcite{cidh2010medidas}, \textcite{basu2010scitotenv} y
  \textcite{etech2010water}. La clasificación es una adaptación propia
  construida sobre \textcite{iso2018iso19011}; la evaluación original no
  clasificó sus resultados como no conformidades.} \\
  \endlastfoot

  % --- 1. CONSULTA ---
  Consulta &
  La operación fue autorizada sin un proceso de consulta previa, libre e informada
  adecuado a la población maya-mam y maya-sipakapense &
  Constatación expresa de la evaluación; pronunciamiento de la {CEACR} de la
  {OIT} (2010) instando a suspender actividades hasta cumplir la consulta &
  Convenio 169 {OIT}, art. 6; {IFC} {PS}7 &
  \NCmayor &
  Detener la adquisición de tierras, la exploración y toda expansión hasta
  consultar debidamente &
  \textbf{No cerrada.} La operación continuó hasta 2017 \\
  \addlinespace

  % --- 2. AMBIENTE (a) ---
  Ambiente &
  Sistema de monitoreo ambiental y línea base insuficientes para atribuir o
  descartar impactos &
  Revisión independiente del {EIA} \parencite{moran2004review}; recomendación de
  la propia evaluación de reforzar el monitoreo &
  {IFC} {PS}1 &
  \NCmayor &
  Reforzar el monitoreo, incluida la participación comunitaria a través de la
  {AMAC} &
  \textbf{Parcial.} Planta de tratamiento de agua (2008) y monitoreo comunitario;
  sin verificación independiente \\
  \addlinespace

  % --- 3. AMBIENTE (b) ---
  Ambiente &
  Riesgo de drenaje ácido de roca por la mineralogía piritosa del yacimiento y del
  desmonte &
  Composición mineralógica declarada en el informe técnico
  \parencite{sec2011ni43101}; advertencia de la evaluación &
  {IFC} {PS}1, {PS}3 &
  \NCmenor &
  Incorporar el drenaje ácido a la planificación del cierre y al monitoreo de
  largo plazo &
  \textbf{No verificable} con la documentación pública disponible \\
  \addlinespace

  % --- 4. AMBIENTE (c) ---
  Ambiente &
  Afectación efectiva de la calidad del agua atribuible a la operación &
  Evidencia \textbf{contradictoria}: metales elevados en población cercana
  \parencite{basu2010scitotenv}, pero por debajo de los valores asociados a
  daño clínico; efectos adversos que «pueden haber comenzado ya»
  \parencite{etech2010water}, frente a monitoreos que sostienen conformidad &
  --- &
  \OBS &
  Línea base pública, cadena de custodia y laboratorio acreditado con validación
  de terceros &
  \textbf{No resuelta.} Ver sección \ref{sec:critico} \\
  \addlinespace

  % --- 5. TRABAJO (a) ---
  Trabajo &
  La empresa habría desalentado activamente la formación de un sindicato &
  Constatación de la evaluación &
  {PIDCP}, art. 22; convenios fundamentales {OIT} &
  \NCmayor &
  Garantizar la libertad sindical y de negociación colectiva &
  \textbf{No verificable} \\
  \addlinespace

  % --- 6. TRABAJO (b) ---
  Trabajo &
  Deficiencias en seguridad y salud ocupacional &
  Constatación de la evaluación &
  {PIDESC}, art. 7 &
  \NCmenor &
  Fortalecer el sistema de gestión de seguridad y salud en el trabajo &
  \textbf{No verificable} \\
  \addlinespace

  % --- 7. TIERRAS ---
  Adquisición de tierras &
  Deficiencias de equidad y transparencia en el proceso de adquisición, y
  afectación de derechos colectivos &
  Constatación de la evaluación &
  {IFC} {PS}5, {PS}7 &
  \NCmenor &
  Revisar el procedimiento de adquisición y reconocer la titularidad colectiva &
  \textbf{No verificable} \\
  \addlinespace

  % --- 8. INVERSIÓN ---
  Inversión económica y social &
  Mitigación insuficiente de los impactos negativos, pese a los aportes efectivos
  en empleo e ingresos &
  Constatación de la evaluación; datos de empleo y producción (sección
  \ref{sec:unidad}) &
  {IFC} {PS}1 &
  \NCmenor &
  Reequilibrar la inversión social hacia la mitigación y no solo hacia el aporte &
  \textbf{Parcial} \\
  \addlinespace

  % --- 9. SEGURIDAD ---
  Seguridad &
  Riesgo para la seguridad personal de los residentes por la presencia de
  seguridad privada armada y de fuerzas públicas &
  Constatación de la evaluación; medidas cautelares {MC} 260-07 de la {CIDH}
  \parencite{cidh2010medidas} &
  {VPSHR}; {IFC} {PS}4 &
  \NCmayor &
  Adoptar formalmente los Principios Voluntarios en Seguridad y Derechos Humanos &
  \textbf{No verificable} la adopción formal \\
  \addlinespace

  % --- 10. REMEDIACIÓN ---
  Acceso a remediación &
  Mecanismos de queja y reclamo débiles &
  Constatación de la evaluación; escalada del caso a la {CAO}/{IFC} y a la
  {CIDH} \parencite{cao2006marlin} &
  {IFC} {PS}1; marco Ruggie, pilar «Remediar» &
  \NCmayor &
  Implantar un mecanismo de quejas conforme a los criterios de eficacia &
  \textbf{No verificable} \\

\end{longtable}
\end{footnotesize}
\end{landscape}

La lectura de la matriz arroja tres conclusiones inmediatas. La primera es que las
no conformidades mayores no se distribuyen de manera uniforme: se concentran en
consulta, monitoreo, libertad sindical, seguridad y remediación, es decir, en los
ejes donde están comprometidos derechos fundamentales o donde falla el sistema en
su conjunto. La segunda es que el eje ambiental, pese a ser el que más atención
pública concentró, no genera la no conformidad más grave por afectación
demostrada, sino por \textbf{insuficiencia del dispositivo de verificación}: no se
sabe lo que ocurrió porque no se midió como debía medirse. La tercera es la más
incómoda y se desarrolla en la sección \ref{sec:sintesis}: la columna de estado de
cierre está dominada por la expresión «no verificable».

% ---------------------------------------------------------------------
\subsection{Hallazgo central 1: ausencia de consulta previa, libre e informada}

Es el hallazgo de mayor gravedad y el que articula todo el caso. La operación fue
autorizada y construida sin que mediara un proceso de consulta previa, libre e
informada adecuado a los pueblos maya-mam y maya-sipakapense que habitan el
territorio afectado \parencite{oncommonground2010hria}.

El criterio incumplido es el artículo 6 del Convenio 169 de la Organización
Internacional del Trabajo, que obliga a consultar a los pueblos interesados
mediante procedimientos apropiados y a través de sus instituciones
representativas, de buena fe y con la finalidad de llegar a un acuerdo
\parencite{oit1989convenio169}. A ello se suma el Performance Standard 7 de la
Corporación Financiera Internacional, exigible por vía contractual en virtud del
préstamo descrito en la sección \ref{sec:unidad} \parencite{ifcsfmarlin}.

La calificación como no conformidad mayor se sostiene en tres razones. La primera
es que el derecho vulnerado es un derecho colectivo de un pueblo indígena,
reconocido en un tratado internacional. La segunda es su carácter \textbf{no
subsanable en el momento en que se detecta}: una consulta es previa o no es
consulta. Una vez otorgada la licencia, construida la mina e iniciada la
producción, no existe acción correctiva capaz de restituir la situación
original; solo caben medidas de reparación, que son otra cosa. La tercera es su
condición de causa raíz: la ausencia de consulta explica el déficit de
legitimidad que arrastró la operación durante toda su vida útil y, con él, la
intensidad del conflicto posterior.

La evaluación recomendó detener la adquisición de tierras, la exploración y
cualquier expansión hasta consultar debidamente. La recomendación fue respaldada
externamente ese mismo año: la Comisión de Expertos en Aplicación de Convenios y
Recomendaciones de la Organización Internacional del Trabajo instó en 2010 a
suspender las actividades hasta dar cumplimiento a la obligación de consulta
\parencite{oncommonground2010hria}.

Conviene registrar también el antecedente. En 2005 el municipio de Sipacapa había
celebrado una consulta comunitaria de buena fe, ejercicio de participación
promovido por la propia población ante la ausencia de un procedimiento estatal.
Su resultado fue mayoritariamente contrario a la actividad minera. Que una
comunidad organice por su cuenta el procedimiento que el Estado no convocó es, en
términos de auditoría, evidencia directa de la magnitud del vacío que se está
constatando.

El estado de cierre de esta no conformidad es inequívoco: la operación continuó
hasta su agotamiento natural en 2017. La no conformidad más grave del expediente
nunca se cerró.

% ---------------------------------------------------------------------
\subsection{Hallazgo central 2: insuficiencia del monitoreo y conflicto de evidencia sobre el agua}

Este hallazgo exige una distinción que la matriz ya anticipa y que constituye el
aporte analítico central de esta sección: \textbf{no es lo mismo afirmar que hubo
contaminación que afirmar que el sistema no permitía saberlo}. Lo primero no está
demostrado; lo segundo sí.

\subsubsection{Lo que sí puede afirmarse}

La revisión independiente del estudio de impacto ambiental elaborada por
\textcite{moran2004review} advirtió deficiencias en el instrumento predictivo
antes incluso del inicio de la producción. La propia evaluación de derechos
humanos recomendó reforzar el monitoreo, incluida la participación comunitaria a
través de la Asociación de Monitoreo Ambiental Comunitario, y mejorar la
planificación del cierre \parencite{oncommonground2010hria}.

La evidencia más concluyente sobre este punto procede, sin embargo, de la
evaluación técnica independiente elaborada por \textcite{etech2010water}, que
comparó las predicciones del estudio de impacto ambiental con las condiciones
efectivamente registradas durante la operación. Sus constataciones son
directamente verificables y cuantificadas:

\begin{itemize}
  \item El período de línea base de calidad de agua previo al inicio de la
        minería fue de \textbf{solo ocho o nueve meses}, insuficiente para
        capturar la variación estacional e interanual.
  \item Para agua subterránea \textbf{se muestrearon únicamente dos
        manantiales}, y no se muestreó en absoluto el acuífero profundo.
  \item No se dispuso de información suficiente sobre niveles y direcciones de
        flujo del agua subterránea, lo que —en palabras del propio informe—
        hace imposible conocer el potencial de migración de contaminantes
        desde las fuentes mineras hacia los receptores.
  \item El cuerpo principal del estudio de impacto ambiental \textbf{no
        incorporó los ensayos geoquímicos}: predijo que el potencial de
        generación de acidez sería bajo, pero sin aportar las tablas ni las
        figuras que lo sustentaran. La evaluación independiente concluyó, por
        el contrario, que los residuos mineros presentan un potencial
        \textbf{moderado a alto} de generar acidez y lixiviar contaminantes.
\end{itemize}

Estas constataciones confirman el riesgo geoquímico que la mineralogía piritosa
descrita en la sección \ref{sec:unidad} hacía previsible
\parencite{sec2011ni43101}.

El incumplimiento, por tanto, no consiste en haber contaminado —extremo no
probado— sino en no haber dispuesto de un sistema de monitoreo y de una línea base
capaces de acreditar la ausencia de contaminación cuando esta fue cuestionada.
Esa carencia se clasifica como no conformidad mayor por su carácter sistémico:
inutiliza el mecanismo de verificación del que depende todo el sistema de gestión
ambiental. Y no es una deficiencia deducida por este informe, sino documentada
por una evaluación técnica independiente.

\subsubsection{Lo que no puede afirmarse}

La evidencia sobre afectación efectiva es contradictoria y debe presentarse como
tal.

Del lado de la afectación, el estudio publicado en \textit{Science of the Total
Environment} por \textcite{basu2010scitotenv}, con participación de la
Universidad de Michigan y de Physicians for Human Rights —organización que
publicó además el informe institucional del mismo trabajo
\parencite{basu2010metales}—, halló que quienes
residían más cerca de la mina presentaban niveles significativamente más altos
de determinados metales —mercurio, cobre, arsénico y zinc en orina— que el grupo
de comparación. Por su parte, \textcite{etech2010water} concluyó que los
efectos adversos sobre el ambiente podían haber comenzado ya como resultado de
las operaciones mineras.

Ahora bien, el propio estudio epidemiológico incorporó dos advertencias que sería
tendencioso omitir. La primera es que \textbf{la mayoría de los metales se
detectaron en concentraciones inferiores a los valores asociados a daño
clínico}. La segunda es que sus autores concluyeron que hacían falta estudios
más robustos antes de determinar si esas elevaciones representaban una amenaza
significativa para la salud, si bien advirtieron que los efectos de la
exposición simultánea a varios elementos podrían ser mayores de lo esperado por
interacción sinérgica. Citar la primera mitad del hallazgo sin la segunda
constituiría una manipulación de la fuente, y este informe no la comete.

Del lado contrario, la empresa y determinados monitoreos oficiales sostuvieron que
no existía contaminación atribuible a la operación.

\subsubsection{Por qué la evidencia no converge}

La discrepancia no obedece necesariamente a mala fe de ninguna de las partes, sino
a diferencias metodológicas que la hacen previsible: distintos puntos de muestreo,
distintas matrices analizadas —orina y sangre frente a agua superficial—,
distintos períodos, distintos laboratorios y, sobre todo, la ausencia de una línea
base pública y anterior al inicio de la operación que permitiera establecer un
término de comparación aceptado por todos.

Un estudio de biomonitoreo humano acredita exposición, no fuente: detectar
arsénico en orina no identifica de dónde procede ese arsénico, que puede tener
origen geogénico en un terreno mineralizado. Precisamente por eso la línea base
previa resulta indispensable, y precisamente por eso su ausencia se clasifica como
no conformidad mayor mientras la afectación misma se registra como observación. La
consecuencia práctica es que el conflicto de evidencia era, en el diseño del
sistema, técnicamente irresoluble desde el principio. Este punto se retoma en la
sección \ref{sec:critico}.

% ---------------------------------------------------------------------
\subsection{Hallazgo central 3: seguridad privada armada y riesgo para las personas}

La evaluación identificó un riesgo para la seguridad personal de los residentes
derivado de la presencia de seguridad privada armada y de fuerzas de seguridad
pública en el entorno de la operación, y recomendó la adopción formal de los
Principios Voluntarios en Seguridad y Derechos Humanos
\parencite{oncommonground2010hria, vpshrsfprincipios}.

Se clasifica como no conformidad mayor porque el bien jurídico comprometido es la
integridad física de las personas, y porque el riesgo no se localiza en un
incidente aislado sino en la configuración misma del dispositivo de seguridad.

Este hallazgo obtuvo la validación externa más contundente de todo el expediente.
El 20 de mayo de 2010 la Comisión Interamericana de Derechos Humanos otorgó
medidas cautelares a favor de los miembros de dieciocho comunidades del pueblo
maya —sipakapense y mam— del departamento de San Marcos, solicitando inicialmente
la suspensión de la mina \parencite{cidh2010medidas}. Que un órgano del sistema
interamericano adopte medidas cautelares implica la apreciación de una situación
de gravedad y urgencia con riesgo de daño irreparable a personas determinadas. Es
el grado más alto de corroboración externa que un hallazgo de auditoría puede
recibir.

En diciembre de 2011 la Comisión modificó las medidas y retiró la solicitud de
suspensión, reorientándolas hacia la garantía de acceso a agua apta para el
consumo humano. La empresa hizo pública esa modificación, subrayando que las
medidas revisadas no requerían suspender las operaciones y que estas habían
continuado con normalidad a lo largo de todo el procedimiento ante la Comisión.

La afirmación es exacta y por eso se recoge. Pero merece una lectura que se
desarrolla en la sección \ref{sec:critico}: durante el año y medio en que
estuvieron vigentes unas medidas cautelares internacionales que pedían la
suspensión de la actividad, la actividad no se suspendió. Ese dato mide, mejor que
ninguna declaración, la capacidad real de estos mecanismos para alterar la
conducta de una operación en curso.

Este eje es, además, el que con mayor nitidez muestra por qué la evaluación
analizada no podía ser una auditoría ambiental convencional. Ninguna cláusula de
la ISO 14001 se pronuncia sobre el porte de armas por personal de seguridad
privada en el entorno de una comunidad indígena.

% ---------------------------------------------------------------------
\subsection{Hallazgo central 4: debilidad del mecanismo de acceso a remediación}

El último hallazgo central es, en apariencia, el más técnico y el menos
llamativo: los mecanismos de queja y reclamo de la operación eran débiles
\parencite{oncommonground2010hria}. Se clasifica, sin embargo, como no
conformidad mayor, y conviene justificar por qué.

Un mecanismo de quejas es el dispositivo mediante el cual una organización detecta
y corrige internamente sus propios fallos. Si funciona, los problemas se resuelven
en el nivel en que se originan. Si no funciona, ninguno de los demás hallazgos
dispone de vía interna de corrección, y el conflicto migra necesariamente hacia
instancias externas. Se trata, en la terminología de la sección \ref{sec:marco},
de un fallo sistémico en sentido estricto: no compromete un requisito, compromete
la capacidad de la organización para cumplir todos los demás.

La trayectoria del caso confirma este diagnóstico de manera casi didáctica. La
queja no se resolvió en la mina: fue presentada por Madre Selva ante la Compliance
Advisor Ombudsman de la Corporación Financiera Internacional en 2005
\parencite{cao2006marlin}. Tampoco se resolvió allí: escaló hasta la Comisión
Interamericana de Derechos Humanos en 2010 \parencite{cidh2010medidas}. Y en
paralelo alcanzó los órganos de control de la Organización Internacional del
Trabajo. Una controversia que en un sistema con mecanismos de reclamo eficaces se
habría tramitado localmente terminó recorriendo tres instancias internacionales a
lo largo de una década.

El criterio incumplido es el Performance Standard 1 de la Corporación Financiera
Internacional y, en el plano conceptual, el tercer pilar del marco «Proteger,
Respetar, Remediar», expuesto en la sección \ref{sec:marco}. La propuesta
correctiva correspondiente se desarrolla en la sección \ref{sec:plan}.

% ---------------------------------------------------------------------
\subsection{Síntesis y respuesta de la empresa}
\label{sec:sintesis}

De los diez hallazgos identificados, cinco se clasifican como no conformidades
mayores, cuatro como menores y uno como observación. Los ejes que concentran la
mayor gravedad son consulta, monitoreo ambiental, libertad sindical, seguridad y
acceso a remediación.

La empresa se comprometió a emitir una respuesta pública y un plan de acción
detallado frente a las recomendaciones, con actualizaciones publicadas en 2010,
2011 y 2017, e integró los aprendizajes del proceso a sus políticas corporativas
globales \parencite{oncommonground2010hria, goldcorp2023marlin}. La operación
cerró en mayo de 2017 con un plan de recuperación ambiental que incluyó la
destrucción de cianuro y el relleno del tajo.

Ahora bien, el elemento más significativo de esta síntesis no es lo que la matriz
afirma, sino lo que no puede afirmar. La columna de estado de cierre está dominada
por la expresión «no verificable»: existe constancia de que la empresa publicó
actualizaciones de su plan de acción, pero no consta una \textbf{verificación
independiente} de que cada recomendación se implementara efectivamente ni de que
cada no conformidad se cerrara.

Esta circunstancia enlaza directamente con la debilidad identificada en la sección
\ref{sec:auditoria} al contrastar el proceso con las fases de la ISO 19011:2018.
La norma concibe la auditoría como un ciclo que se cierra verificando la
implementación de las acciones correctivas; aquí esa verificación quedó
encomendada a la propia organización auditada, sin que el equipo auditor
conservara mandato para comprobarla. El resultado es que, siete años después del
cierre de la operación y quince después de la publicación del informe, no es
posible determinar con la documentación pública disponible cuántas de las
recomendaciones se cumplieron.

En la lógica de una auditoría, una no conformidad cuyo cierre no puede verificarse
permanece abierta. Esa es la conclusión con la que esta sección entrega el caso al
análisis crítico de la sección \ref{sec:critico}.
